\documentclass[conference]{IEEEtran}
\usepackage{cite}
\usepackage{amsmath,amssymb,amsfonts}
\usepackage{graphicx}
\usepackage{booktabs}
\usepackage{url}

\graphicspath{{../data/figures/}}

\begin{document}

\title{Closed-Loop versus Decomposed Go-to-Goal Control of a Unicycle Robot
Under Actuation Noise:\\ A Simulation-Based Comparison}

\author{\IEEEauthorblockN{Mohamed Tawakol}
\IEEEauthorblockA{\textit{MSc Robotics Engineering -- Research Track II}\\
\textit{University of Genoa}\\
Genoa, Italy\\
7952397@studenti.unige.it}}

\maketitle

\begin{abstract}
Reaching a goal position is the most elementary navigation capability of a
mobile robot, yet the choice of the low-level go-to-goal strategy affects both
efficiency and robustness. This paper compares two widely used strategies for a
kinematic unicycle robot: a \emph{proportional closed-loop} controller, which
continuously and simultaneously regulates the linear and angular velocities as
a function of the distance and bearing to the goal, and a \emph{decomposed
rotate-then-translate} baseline, which first aligns the heading and then
drives straight, re-aligning whenever the heading error exceeds a threshold.
The two controllers are evaluated in a ROS~2-based simulation under three
levels of additive Gaussian actuation noise, with a $2\times3$ full factorial
design and $N=25$ paired trials per condition. The closed-loop controller
reaches the goal on average 32\% faster at every noise level (paired
$t$-test, $p<0.001$, $|d|>3$), while the baseline yields straighter paths
(efficiency $0.99$ vs.\ $0.93$ in the noise-free case) but degrades roughly
twice as fast as noise increases. Contrary to our third hypothesis, the
time-to-goal advantage of the closed-loop law does not grow significantly
with the noise level ($p=0.54$). All experiments, data, and figures are fully
reproducible through the accompanying Jupyter notebook, which interacts with
the ROS~2 experiment node.
\end{abstract}

\begin{IEEEkeywords}
mobile robots, unicycle model, go-to-goal control, actuation noise,
experimental methodology, ROS~2
\end{IEEEkeywords}

\section{Introduction and Research Question}
Point-to-point navigation --- driving a wheeled robot from its current pose to
a goal position --- underlies virtually every mobile-robot application, from
service robotics to warehouse logistics. When a global planner provides a
sequence of waypoints, a low-level \emph{go-to-goal} controller is responsible
for actually reaching each of them. Two families of strategies are commonly
taught and deployed. The first \emph{decomposes} the motion: the robot rotates
in place until it points towards the goal and then translates along a straight
line; the behaviour is simple to implement and to predict, and the resulting
path is (nominally) the shortest possible one. The second family uses
\emph{closed-loop} feedback laws that update the linear and angular
velocities simultaneously and continuously, typically as proportional
functions of the polar coordinates of the goal in the robot frame
\cite{aicardi1995closed,siegwart2011introduction}. Since the two velocity
channels are never idle at the same time, the robot keeps advancing while it
corrects its heading, at the price of a curved trajectory.

While the closed-loop approach is generally considered superior, the
decomposed strategy is still widespread in educational settings and in simple
industrial deployments, where its straight-line paths are seen as an
advantage --- for instance in narrow aisles, where lateral deviation matters
more than travel time. Moreover, real robots execute velocity commands
imperfectly: gear backlash, wheel slip, discretized motor drivers, and
unmodelled dynamics all perturb the executed linear and angular velocities.
Intuition suggests that a strategy that corrects errors continuously should
tolerate such disturbances better than one that only reacts once a threshold
is crossed, but intuition is not evidence. It is therefore legitimate --- and,
to our knowledge, not yet answered quantitatively in the literature --- to ask
whether the theoretical advantages of continuous feedback survive, or even
grow, when actuation becomes unreliable.

Beyond its immediate object, the study is also designed as an exercise in
sound experimental methodology: hypotheses are stated before data collection,
all variables are explicitly classified as independent, dependent, or
controlled, trials are paired through shared random seeds, and every number
and figure in this paper can be regenerated by executing a single Jupyter
notebook that drives the ROS~2 experiment node.

Following the compressed PICO structure \cite{bonsignorio2015toward}, the
research question addressed in this work is:

\begin{quote}
\emph{Does a proportional closed-loop go-to-goal controller improve
time-to-goal and path efficiency, compared to a rotate-then-translate
baseline, for a unicycle mobile robot navigating to random goals under
increasing actuation noise?}
\end{quote}

Three falsifiable hypotheses were formulated \emph{before} data collection:
\begin{itemize}
\item \textbf{H1}: the proportional controller achieves a lower time-to-goal
than the baseline at every noise level;
\item \textbf{H2}: the baseline achieves a higher path efficiency in the
noise-free case, but its efficiency degrades faster as noise increases;
\item \textbf{H3}: the time-to-goal advantage of the proportional controller
increases with the noise level (interaction effect).
\end{itemize}

\section{Literature Analysis}
The kinematic control of unicycle-like vehicles is a classical subject.
Brockett's theorem \cite{brockett1983asymptotic} establishes that no smooth,
time-invariant, pure-state feedback can asymptotically stabilize a
nonholonomic vehicle to a full pose, which historically motivated two research
directions: discontinuous or time-varying laws for pose stabilization
\cite{astolfi1999exponential}, and smooth laws for the easier
\emph{position-only} stabilization problem considered here. Within the second
direction, Aicardi \emph{et al.} \cite{aicardi1995closed} introduced the polar
coordinate transformation and the Lyapunov-based proportional law on which our
``proposed'' controller is modelled: expressing the goal through its distance
$\rho$ and bearing error $\alpha$, the simple feedback
$v=k_\rho\,\rho$, $\omega=k_\alpha\,\alpha$ yields global convergence with
smooth, arc-shaped trajectories. This law, popularized by the textbook of
Siegwart \emph{et al.} \cite{siegwart2011introduction}, remains the standard
entry point for go-to-goal control. In the closely related trajectory-tracking
setting, Kanayama \emph{et al.} \cite{kanayama1990stable} proposed a stable
tracking scheme whose linearization is again a proportional feedback on the
posture error, reinforcing the centrality of proportional laws in this
domain. Later refinements focused on \emph{graceful} motion: Park and Kuipers
\cite{park2011smooth} derived a smooth control law with curvature-bounded
trajectories for differential-drive robots, confirming that continuous
feedback produces faster and smoother motion than decomposed schemes, though
without a systematic statistical comparison under actuation noise. A common
trait of these control-theoretic works is that performance is demonstrated by
convergence proofs and a handful of illustrative trajectories, not by
repeated randomized experiments: robustness to stochastic actuation error is
argued qualitatively, when discussed at all.

A second body of literature addresses obstacle-aware local navigation. The
dynamic window approach \cite{fox1997dynamic} and its descendants, up to the
ROS~2 Nav2 stack \cite{macenski2020marathon}, embed goal-seeking behaviour
into velocity-space optimization. These controllers presuppose exactly the
kind of elementary go-to-goal capability studied here, but their complexity
confounds the analysis of the low-level strategy itself; benchmarking papers
about Nav2 typically compare whole planning pipelines rather than the
underlying motion primitives. Interestingly, Nav2 itself ships a
``rotate-to-heading, then drive'' behaviour in its rotation shim controller,
precisely because pure closed-loop controllers may start driving in the
wrong direction when the initial heading error is large --- evidence that the
decomposed strategy is not merely a pedagogical strawman but a design option
that practitioners actively choose.

Finally, a growing literature on experimental methodology in robotics
\cite{bonsignorio2015toward,macenski2022ros2} stresses replicability:
experiments should isolate independent variables, report aggregate statistics
over repeated randomized trials, and publish code and data. Interactive
notebooks are increasingly adopted as the vehicle for such reproducible
workflows \cite{kluyver2016jupyter}.

\textbf{Gap.} Across these three threads, the closed-loop and decomposed
go-to-goal strategies are either assumed (navigation stacks), analysed only
in the noise-free deterministic setting (control-theoretic works), or
compared qualitatively. To the best of our knowledge, no work reports a
controlled, statistically analysed comparison of the two strategies under
increasing \emph{actuation} noise, with paired randomized trials and public
code. This paper fills that small but well-defined gap, and does so with a
fully reproducible ROS~2 + notebook pipeline.

\section{Methodological Proposal}
\subsection{Robot and noise model}
The robot is modelled as a kinematic unicycle with state $(x,y,\theta)$ and
inputs $(v,\omega)$, integrated with explicit Euler at $\Delta t=0.05$\,s:
\begin{equation}
x^+ = x + v\cos\theta\,\Delta t,\quad
y^+ = y + v\sin\theta\,\Delta t,\quad
\theta^+ = \theta + \omega\,\Delta t .
\end{equation}
Actuation imperfection is modelled as additive zero-mean Gaussian noise on the
\emph{executed} velocities at every step:
\begin{equation}
v = v_{\mathrm{cmd}} + \mathcal{N}(0,\sigma\,V_{\max}),\qquad
\omega = \omega_{\mathrm{cmd}} + \mathcal{N}(0,\sigma\,W_{\max}),
\end{equation}
where $\sigma\in\{0.0,\,0.1,\,0.3\}$ is the dimensionless noise level and
$V_{\max}=2$\,m/s, $W_{\max}=2$\,rad/s are the saturation limits. Executed
linear velocities are clipped at zero (no reverse motion).

\subsection{Controllers}
Both controllers observe the goal through its polar coordinates in the robot
frame: the distance $\rho=\|(x_g-x,\,y_g-y)\|$ and the bearing error
$\alpha=\operatorname{atan2}(y_g-y,\,x_g-x)-\theta$, wrapped to
$(-\pi,\pi]$. The \textbf{proportional closed-loop controller} (proposed) is
the classical polar law \cite{aicardi1995closed} with a cosine gating that
slows the robot down when badly oriented:
\begin{equation}
v = k_\rho\,\rho\,\max(\cos\alpha,0),\qquad
\omega = k_\alpha\,\alpha ,
\end{equation}
with $k_\rho=1.2$, $k_\alpha=3.0$. The gating term preserves the convergence
properties of the original law while preventing the robot from driving
forward when the goal is behind it. The \textbf{rotate-then-translate
baseline} implements the decomposed strategy as a two-state machine: it
rotates in place at $\omega=\pm1.5$\,rad/s while $|\alpha|>0.2$\,rad, and
otherwise translates at a constant $v=1$\,m/s with $\omega=0$; whenever noise
pushes the heading error back above the threshold, it stops and re-aligns.
Both outputs are saturated to $[0,V_{\max}]$ and $[-W_{\max},W_{\max}]$. The
gains were fixed once, by hand, so that both controllers use comparable
fractions of the available velocity envelope; neither was tuned per
condition, which keeps the comparison fair but is acknowledged as a
limitation in Sec.~VI.

\subsection{Software architecture}
The experiment is implemented as a ROS~2 (Jazzy) Python package,
\texttt{goal\_nav\_experiment}\footnote{Code, executed notebook, and data:
\url{https://github.com/Tawakoll/RT2_Final_Project}}. A dedicated
\emph{experiment node} receives JSON trial requests on the topic
\texttt{/experiment/run} (controller name, noise level, seed), simulates the
trial, publishes the resulting metrics on \texttt{/experiment/result}, and
appends every trial to a CSV log. A Jupyter notebook acts as the ROS~2 client:
it launches the node, requests single interactive trials through
\texttt{ipywidgets} controls, runs the full experimental campaign, and
produces all statistics and figures reported here. Because every trial is
driven by a seeded random generator, the entire study is exactly
reproducible \cite{kluyver2016jupyter}.

The simulated experiment makes the following explicit assumptions: (i) purely
kinematic motion (no inertia, slip, or actuator dynamics beyond the noise
term); (ii) perfect localization (the controllers read the true pose);
(iii) an obstacle-free $11\times11$\,m arena; (iv) noise that is additive,
Gaussian, zero-mean, and white. These assumptions isolate the independent
variable of interest --- the control strategy --- from confounding factors.

\section{Experiment Design}
\subsection{Variables}
\textbf{Independent variables}: controller type (proportional,
rotate-then-translate) and actuation noise level
($\sigma\in\{0.0,0.1,0.3\}$), in a $2\times3$ full factorial design.
\textbf{Dependent variables}: time-to-goal [s]; path length [m]; path
efficiency, defined as the ratio between the distance actually covered
towards the goal and the travelled path length (1 = straight line); success
rate (goal reached within the timeout). \textbf{Control variables}: arena
size, start pose $(5.5,5.5,0)$, velocity limits, integration step, goal
tolerance ($0.15$\,m), timeout ($60$\,s), controller gains, and the goal
distribution (uniform in $[1,10]^2$, at least $3$\,m from the start).

\subsection{Protocol}
For every condition, $N=25$ trials are executed. The same 25 seeds are reused
across all six conditions, so that both controllers face \emph{exactly} the
same sequence of goals under the same noise realization schedule: the design
is fully \emph{paired}, which increases statistical power and justifies the
use of the paired $t$-test. The number of trials was chosen after a small
pilot (5 runs per condition) showed low variability and a very large effect
size for time-to-goal. In total, $2\times3\times25=150$ trials are executed
automatically by the notebook in a few seconds of wall-clock time.

Confounding factors are handled as follows. Environmental variability is
eliminated by construction (the arena is fixed and empty); goal-position
variability, which strongly affects all metrics, is neutralized by the paired
seeding; the stochastic noise realization is the only residual source of
within-pair variability, and it is averaged out by the repeated trials.
Because the simulation is deterministic given the seed, re-running the
campaign reproduces the dataset bit for bit --- failures, if any, are included
in the analysis rather than discarded.

\subsection{Statistical analysis}
For each metric and noise level, the two controllers are compared with a
two-tailed paired $t$-test at significance level $\alpha=0.05$; the Wilcoxon
signed-rank test is used as a non-parametric robustness check, and Cohen's
$d$ on the paired differences quantifies the effect size. H3 is tested by
comparing the paired time-to-goal differences at $\sigma=0.3$ against those
at $\sigma=0.0$ with a Welch two-sample $t$-test.


\section{Results and Visualization}
% TODO: aggregate metrics, hypothesis tests, qualitative evaluation

\section{Discussion and Limitations}
% TODO

\section{Conclusion}
% TODO

\end{document}
